\chapter{Organizzazione del corso}
\label{cap:organizzazione}

\section{Percorso in tre parti}

\begin{center}\small
\begin{tabular}{lp{4cm}p{7cm}}
\toprule
& Contenuto & Obiettivi di apprendimento \\
\midrule
Parte I & Modellazione & riconoscere un legame fra variabili (attivazione, massimo,
big-M, se e solo se\ldots) e dimostrare che il modello lo impone davvero \\
Parte II & I problemi & applicare i legami a tre famiglie di problemi reali e ai problemi misti,
dal modello al codice Gurobi \\
Parte III & Il corso & mettere alla prova quanto imparato con le domande di
modellazione aggiuntive \\
\bottomrule
\end{tabular}
\end{center}

\section{Il formato di ogni esercizio (e dell'esame)}

Ogni problema del corso --- e ogni domanda d'esame --- segue lo stesso schema
in quattro quesiti:
\begin{enumerate}
\item \textbf{Modello.} Scrivere il MILP: variabili (con il conteggio),
      obiettivo, vincoli con la loro descrizione, e --- se ci sono due
      famiglie di variabili collegate --- spiegare il legame: qual è
      l'implicazione, e come il vincolo (o l'ottimo) la impone, nei due versi.
\item \textbf{Istanza.} Scrivere il modello per un'istanza numerica piccola.
\item \textbf{Euristica.} Progettare un algoritmo costruttivo che trovi una
      soluzione ammissibile dell'istanza, ed eseguirlo passo per passo per
      ottenere un upper bound (lower bound se il problema è di massimo).
\item \textbf{Duale.} Scrivere il duale del rilassamento LP\index{rilassamento lineare}, per il modello
      generale e per l'istanza, e costruire a mano una soluzione duale
      ammissibile per ottenere il bound dal lato opposto.
\end{enumerate}
Il $\lb \le \zmilp \le \ub$ che ne risulta è il filo conduttore del corso: un
modello non si limita a scriverlo, lo si stringe da entrambi i lati prima di
affidarlo al solver.

\section{Le domande aggiuntive e le loro soluzioni}
Ogni problema si chiude con due o tre \emph{domande di modellazione aggiuntive}:
variano un dato o aggiungono un vincolo, e chiedono come cambia il modello e
quanto vale il nuovo ottimo. Sono la parte su cui conviene esercitarsi, perché
il modello base lo si legge, la variante la si scrive.

\textbf{Le soluzioni delle domande aggiuntive sono riservate ai docenti}: stanno
in un fascicolo a parte, che non viene pubblicato. Quello che resta a
disposizione di tutti --- ed è quanto serve per sapere che cosa ci si aspetta ---
è il \emph{sandwich svolto per intero} su una variante di ciascun esercizio dei
capitoli «Assegnamento e scheduling» e «Localizzazione e copertura»: il modello della
variante, un'euristica costruttiva ammissibile eseguita passo per passo, un
certificato duale costruito a mano e la tabella dei bound. È il modello di
risposta che l'esame si aspetta.

\section{Criteri di valutazione}

\begin{center}\small
\begin{tabular}{lr}
\toprule
Dimensione & Peso \\
\midrule
Correttezza del modello (variabili, obiettivo, vincoli) & 35\% \\
Dimostrazione del legame fra le variabili (nei due versi) & 25\% \\
Euristica costruttiva ed esecuzione corretta & 20\% \\
Duale del rilassamento e soluzione duale ammissibile & 20\% \\
\bottomrule
\end{tabular}
\end{center}

\section{Domande tipiche di discussione}
\begin{itemize}
\item Il vincolo di link è aggregato o disaggregato? Che differenza fa sul
      rilassamento LP\index{rilassamento lineare}?
\item Il verso opposto dell'implicazione è imposto dal vincolo o segue
      dall'ottimo? Come lo si dimostra?
\item Qual è il big-M\index{big-M} più piccolo che si può giustificare dai dati?
\item L'euristica trova l'ottimo? Come si fa a saperlo senza il solver?
\item Il duale a mano è ottimo per il rilassamento, o solo ammissibile?
\end{itemize}

\section{Gli errori più comuni}
\begin{enumerate}
\item Scrivere un'implicazione la cui tesi è vera comunque (controinversa
      vacua): prima di scriverla, verificare che antecedente e conseguente
      siano entrambi genuini.
\item Dimostrare un solo verso di un'implicazione «imposta dal vincolo» ---
      servono sempre entrambi gli ``if''.
\item Confondere una relazione «imposta dal vincolo» con una che «segue
      dall'ottimo»: la seconda richiede l'argomento di scambio\index{argomento di scambio} in sei passi,
      non basta dire «si vede che conviene».
\item Concludere «in ogni soluzione ottima» quando il coefficiente è solo
      $\ge 0$ (non $>0$): la conclusione corretta è più debole, «esiste un
      ottimo in cui\ldots».
\item Usare un big-M\index{big-M} enorme «per sicurezza»: peggiora il rilassamento LP\index{rilassamento lineare}
      senza bisogno.
\item Dimenticare che un vincolo di mutua esclusione ($A+B\le 1$) non implica
      che almeno uno dei due valga 1.
\item Scrivere il duale del rilassamento sbagliando il verso o il segno di
      una variabile duale rispetto al verso del vincolo primale.
\item Costruire una soluzione duale ammissibile ma non verificarne
      l'ammissibilità su tutti i vincoli.
\end{enumerate}

\section{Riproducibilità}
\begin{lstlisting}[style=output]
python3 -m pip install gurobipy matplotlib pandas
python3 python/esegui_tutti.py       # rigenera dati, risultati, figure e notebook
python3 python/verifica_numeri.py    # verifica ogni numero citato nella dispensa
\end{lstlisting}
